\documentclass[11pt]{article}

\usepackage[a4paper,margin=1in]{geometry}
\usepackage{amsmath,amssymb,amsthm,mathtools}
\usepackage{mathrsfs}
\usepackage{enumitem}
\usepackage{array}
\usepackage{hyperref}
\usepackage{tikz-cd}

\hypersetup{
	colorlinks=true,
	linkcolor=blue,
	urlcolor=blue,
	citecolor=blue
}

\setlength{\parskip}{0.6em}
\setlength{\parindent}{0pt}

\newtheorem{theorem}{Theorem}[section]
\newtheorem{proposition}[theorem]{Proposition}
\newtheorem{lemma}[theorem]{Lemma}
\newtheorem{corollary}[theorem]{Corollary}
\theoremstyle{definition}
\newtheorem{definition}[theorem]{Definition}
\newtheorem{example}[theorem]{Example}
\newtheorem{exercise}[theorem]{Exercise}
\theoremstyle{remark}
\newtheorem{remark}[theorem]{Remark}

\newcommand{\F}{\mathbb{F}}
\newcommand{\C}{\mathbb{C}}
\newcommand{\Z}{\mathbb{Z}}
\newcommand{\Q}{\mathbb{Q}}
\newcommand{\R}{\mathbb{R}}
\newcommand{\cO}{\mathcal{O}}
\newcommand{\Syz}{\operatorname{Syz}}
\newcommand{\im}{\operatorname{im}}
\newcommand{\kerR}{\operatorname{ker}}
\newcommand{\Spec}{\operatorname{Spec}}
\newcommand{\Hom}{\operatorname{Hom}}
\newcommand{\Ann}{\operatorname{Ann}}
\newcommand{\rank}{\operatorname{rank}}

\title{Lecture Notes on Syzygies\\
	\large From Linear Dependence to Relations in Commutative Algebra}
\author{}
\date{}

\begin{document}
	
	\maketitle
	
	\tableofcontents
	
	\section{Motivation from linear algebra}
	
	In linear algebra, linear dependence of vectors \(\textbf{v}_1,\dots,\textbf{v}_r\in V\) is expressed by an equation
	\[
	c_1\textbf{v}_1+\cdots+c_r\textbf{v}_r=\textbf{0}
	\]
	for some scalars \(c_1,\dots,c_r\in\F\).
	
	In commutative algebra, one studies an analogous notion for elements of a ring, but now the coefficients are allowed to lie in the ring \(R\) itself.
	
	This is the basic idea behind the notion of a \emph{syzygy}.
	
	\subsection*{Guiding question}
	
	Suppose \(f_1,\dots,f_r\in R\).  
	Can we find elements \(a_1,\dots,a_r\in R\) such that
	\[
	a_1f_1+\cdots+a_rf_r=0?
	\]
	If yes, then the tuple \((a_1,\dots,a_r)\) records a relation among the \(f_i\).  
	That relation is called a syzygy.
	
	\section{From linear dependence to syzygies}
	
	Let us compare the two situations.
	
	\subsection{Linear algebra}
	
	If \(V\) is a vector space over \(\F\), then a relation among vectors
	\[
	\textbf{v}_1,\dots,\textbf{v}_r\in V
	\]
	has the form
	\[
	c_1\textbf{v}_1+\cdots+c_r\textbf{v}_r=\textbf{0},
	\qquad c_i\in \F.
	\]
	The coefficients come from the field \(\F\).
	
	\subsection{Commutative algebra}
	
	If \(R\) is a commutative ring and
	\[
	f_1,\dots,f_r\in R,
	\]
	then a relation among these elements has the form
	\[
	a_1f_1+\cdots+a_rf_r=0,
	\qquad a_i\in R.
	\]
	The coefficients now come from the ring \(R\) itself.
	
	\begin{remark}
		This is richer than linear algebra because coefficients are no longer just constants; they may be polynomials or other ring elements.
	\end{remark}
	
	\section{Definition of syzygy}
	
	\begin{definition}
		Let \(R\) be a commutative ring, and let \(f_1,\dots,f_r\in R\).
		A \emph{syzygy} among \(f_1,\dots,f_r\) is a tuple
		\[
		(a_1,\dots,a_r)\in R^r
		\]
		such that
		\[
		a_1f_1+\cdots+a_rf_r=0.
		\]
	\end{definition}
	
	\begin{definition}
		The set of all syzygies among \(f_1,\dots,f_r\) is called the \emph{syzygy module} of \(f_1,\dots,f_r\), denoted
		\[
		\Syz(f_1,\dots,f_r)
		=
		\left\{
		(a_1,\dots,a_r)\in R^r \;\middle|\; a_1f_1+\cdots+a_rf_r=0
		\right\}.
		\]
	\end{definition}
	
	\begin{proposition}
		\(\Syz(f_1,\dots,f_r)\) is an \(R\)-submodule of \(R^r\).
	\end{proposition}
	
	\begin{proof}
		Let \((a_1,\dots,a_r)\) and \((b_1,\dots,b_r)\) be syzygies. Then
		\[
		a_1f_1+\cdots+a_rf_r=0,
		\qquad
		b_1f_1+\cdots+b_rf_r=0.
		\]
		Adding these equations gives
		\[
		(a_1+b_1)f_1+\cdots+(a_r+b_r)f_r=0,
		\]
		so the sum is again a syzygy.
		
		If \(c\in R\), then multiplying the first equation by \(c\) gives
		\[
		(ca_1)f_1+\cdots+(ca_r)f_r=0.
		\]
		So scalar multiples are again syzygies.
	\end{proof}
	
	\section{The first basic example}
	
	Let
	\[
	R=\C[x,y].
	\]
	Choose
	\[
	f_1=x,\qquad f_2=y.
	\]
	A syzygy of \(f_1,f_2\) is a pair \((a_1,a_2)\in R^2\) such that
	\[
	a_1x+a_2y=0.
	\]
	
	\begin{example}
		The pair
		\[
		(-y,x)
		\]
		is a syzygy of \(x\) and \(y\), because
		\[
		(-y)x+xy=-xy+xy=0.
		\]
	\end{example}
	
	This is the first example one should remember.
	
	\subsection*{Interpretation}
	
	The pair \((-y,x)\) means that the generators \(x\) and \(y\) are not independent. There is a precise way to combine them so that they cancel:
	\[
	(-y)\cdot x + x\cdot y = 0.
	\]
	
	\section{Why a syzygy is a tuple in \(R^r\)}
	
	Suppose we have chosen elements
	\[
	f_1,\dots,f_r\in R.
	\]
	A relation among them must specify one coefficient for each \(f_i\). Therefore the data of a relation is naturally a tuple
	\[
	(a_1,\dots,a_r)\in R^r.
	\]
	The equation
	\[
	a_1f_1+\cdots+a_rf_r=0
	\]
	says that this tuple produces the zero relation.
	
	So:
	
	\begin{itemize}[leftmargin=2em]
		\item \(R^r\) is the space of all coefficient choices,
		\item a syzygy is a coefficient choice that gives zero.
	\end{itemize}
	
	\section{Syzygies as kernels}
	
	This is the cleanest conceptual viewpoint.
	
	Given \(f_1,\dots,f_r\in R\), define an \(R\)-linear map
	\[
	\varphi:R^r\to R
	\]
	by
	\[
	\varphi(a_1,\dots,a_r)=a_1f_1+\cdots+a_rf_r.
	\]
	
	Then by definition,
	\[
	\Syz(f_1,\dots,f_r)=\kerR(\varphi).
	\]
	
	Thus a syzygy is simply an element of the kernel of the map determined by the chosen generators.
	
	\begin{remark}
		This is the most penetrating viewpoint:
		\[
		\text{syzygy}=\text{kernel element of the generator map}.
		\]
	\end{remark}
	
	\subsection{The example \(x,y\)}
	
	Let
	\[
	\varphi:\C[x,y]^2\to \C[x,y],
	\qquad
	\varphi(a_1,a_2)=a_1x+a_2y.
	\]
	Then
	\[
	\Syz(x,y)=\kerR(\varphi).
	\]
	Since
	\[
	\varphi(-y,x)=(-y)x+xy=0,
	\]
	the pair \((-y,x)\) lies in the kernel.
	
	\section{Relation with ideals}
	
	If \(f_1,\dots,f_r\in R\), they generate the ideal
	\[
	(f_1,\dots,f_r)
	=
	\{a_1f_1+\cdots+a_rf_r\mid a_i\in R\}.
	\]
	So there is a natural surjective map
	\[
	\pi:R^r\to (f_1,\dots,f_r),
	\qquad
	(a_1,\dots,a_r)\mapsto a_1f_1+\cdots+a_rf_r.
	\]
	
	The kernel of this map is the module of syzygies:
	\[
	0\to \Syz(f_1,\dots,f_r)\to R^r\xrightarrow{\pi}(f_1,\dots,f_r)\to 0.
	\]
	
	This exact sequence says:
	
	\begin{itemize}[leftmargin=2em]
		\item \(R^r\) remembers the chosen generators,
		\item the image is the ideal they generate,
		\item the kernel records the hidden relations among them.
	\end{itemize}
	
	\section{All syzygies of \(x\) and \(y\)}
	
	We now study the example
	\[
	R=k[x,y],\qquad f_1=x,\qquad f_2=y
	\]
	more carefully.
	
	Suppose
	\[
	ax+by=0
	\]
	for some \(a,b\in R\). We claim that every such relation is a multiple of \((-y,x)\).
	
	\begin{proposition}
		In \(R=k[x,y]\),
		\[
		\Syz(x,y)=R\cdot(-y,x).
		\]
		That is, every syzygy of \(x\) and \(y\) has the form
		\[
		(-cy,cx)
		\]
		for some \(c\in R\).
	\end{proposition}
	
	\begin{proof}
		Suppose
		\[
		ax+by=0.
		\]
		Then
		\[
		ax=-by.
		\]
		Hence \(ax\) is divisible by \(y\). Since \(x\) and \(y\) are relatively prime in \(k[x,y]\), it follows that \(a\) must be divisible by \(y\). So we may write
		\[
		a=-cy
		\]
		for some \(c\in R\). Substituting gives
		\[
		(-cy)x+by=0,
		\]
		hence
		\[
		y(-cx+b)=0.
		\]
		Since \(R\) is an integral domain,
		\[
		-cx+b=0,
		\]
		so
		\[
		b=cx.
		\]
		Therefore
		\[
		(a,b)=(-cy,cx)=c(-y,x).
		\]
	\end{proof}
	
	\begin{remark}
		This is a beautiful example: there are infinitely many syzygies, but they are all generated by one basic syzygy.
	\end{remark}
	
	\section{A second example}
	
	Let
	\[
	R=k[x,y]
	\]
	and consider
	\[
	f_1=x^2,\qquad f_2=xy,\qquad f_3=y^2.
	\]
	We look for triples \((a,b,c)\in R^3\) such that
	\[
	ax^2+bxy+cy^2=0.
	\]
	
	Two obvious syzygies are
	\[
	(-y,x,0)
	\]
	and
	\[
	(0,-y,x).
	\]
	
	Indeed,
	\[
	(-y)x^2+x(xy)+0\cdot y^2=-x^2y+x^2y=0,
	\]
	and
	\[
	0\cdot x^2+(-y)(xy)+x(y^2)=-xy^2+xy^2=0.
	\]
	
	\begin{remark}
		This shows how syzygies arise from different ways of forming the same product.
	\end{remark}
	
	\section{Syzygies and presentations}
	
	Suppose \(M\) is an \(R\)-module generated by elements
	\[
	m_1,\dots,m_r.
	\]
	Then there is a surjective map
	\[
	\pi:R^r\to M,
	\qquad
	(a_1,\dots,a_r)\mapsto a_1m_1+\cdots+a_rm_r.
	\]
	Its kernel is the module of relations among the generators.
	
	\begin{definition}
		The kernel of the map \(R^r\to M\) is called the \emph{module of first syzygies} of the chosen generators of \(M\).
	\end{definition}
	
	Thus one has an exact sequence
	\[
	0\to K_1\to R^r\to M\to 0,
	\]
	where
	\[
	K_1=\kerR(R^r\to M).
	\]
	
	This is called a \emph{presentation} of \(M\): the module \(M\) is described by generators and relations.
	
	\section{Higher syzygies}
	
	The first syzygy module may itself have generators and relations.
	
	Suppose
	\[
	0\to K_1\to F_0\to M\to 0
	\]
	with \(F_0\) free. If \(K_1\) is generated by some elements, then there is a surjection
	\[
	F_1\to K_1\to 0
	\]
	from another free module \(F_1\). The kernel of \(F_1\to K_1\) consists of relations among the first syzygies.
	
	These are called \emph{second syzygies}.
	
	Continuing in this way, one gets
	\[
	\cdots \to F_2\to F_1\to F_0\to M\to 0.
	\]
	
	\begin{definition}
		An exact sequence
		\[
		\cdots \to F_2\to F_1\to F_0\to M\to 0
		\]
		with each \(F_i\) a free \(R\)-module is called a \emph{free resolution} of \(M\).
	\end{definition}
	
	Interpretation:
	\[
	\begin{array}{c|c}
		\text{Module} & \text{Meaning} \\
		\hline
		F_0 & \text{generators of }M \\
		F_1 & \text{first syzygies} \\
		F_2 & \text{second syzygies} \\
		F_3 & \text{third syzygies} \\
		\vdots & \vdots
	\end{array}
	\]
	
	\section{A concrete exact sequence}
	
	For the ideal
	\[
	I=(x,y)\subseteq k[x,y],
	\]
	consider the map
	\[
	\varphi:R^2\to I,
	\qquad
	\varphi(a,b)=ax+by.
	\]
	This map is surjective, and its kernel is generated by \((-y,x)\). Therefore we have an exact sequence
	\[
	0\to R \xrightarrow{\psi} R^2 \xrightarrow{\varphi} I\to 0,
	\]
	where
	\[
	\psi(c)=(-cy,cx).
	\]
	
	This says:
	
	\begin{itemize}[leftmargin=2em]
		\item the ideal \(I\) has two generators \(x,y\),
		\item these generators satisfy one basic relation,
		\item all other relations come from this one.
	\end{itemize}
	
	\section{Comparison with linear algebra}
	
	It is worth returning to the analogy with linear algebra.
	
	\subsection{Vector spaces}
	
	A map
	\[
	T:\F^r\to V,
	\qquad
	(c_1,\dots,c_r)\mapsto c_1\textbf{v}_1+\cdots+c_r\textbf{v}_r
	\]
	has kernel equal to the space of linear relations among \(\textbf{v}_1,\dots,\textbf{v}_r\).
	
	\subsection{Modules over rings}
	
	A map
	\[
	\varphi:R^r\to M,
	\qquad
	(a_1,\dots,a_r)\mapsto a_1m_1+\cdots+a_rm_r
	\]
	has kernel equal to the module of syzygies among \(m_1,\dots,m_r\).
	
	So syzygies are the natural generalization of linear dependence from vector spaces to modules.
	
	\section{Why syzygies matter}
	
	Syzygies matter because they reveal hidden structure.
	
	Knowing generators of an ideal or module is only the first step. One also wants to know:
	
	\begin{itemize}[leftmargin=2em]
		\item how the generators depend on one another,
		\item how complicated those dependencies are,
		\item and whether there are higher layers of dependencies.
	\end{itemize}
	
	Thus syzygies are a way of measuring structure beyond generators.
	
	\subsection*{Slogan}
	\[
	\text{generators} \longrightarrow \text{relations} \longrightarrow \text{relations among relations} \longrightarrow \cdots
	\]
	
	\section{A brief geometric remark}
	
	In algebraic geometry, one often studies a variety by the equations defining it inside projective space. Those equations may satisfy relations among themselves. These relations are syzygies.
	
	So:
	
	\begin{itemize}[leftmargin=2em]
		\item equations describe the variety,
		\item syzygies describe how the equations fit together.
	\end{itemize}
	
	This is why syzygies are central in modern algebraic geometry.
	
	\section{Exercises}
	
	\begin{exercise}
		Let \(R=\C[x,y]\). Show directly that
		\[
		(-y,x)
		\]
		is a syzygy of \(x\) and \(y\).
	\end{exercise}
	
	\begin{exercise}
		Let \(R=k[x,y]\). Prove that every syzygy of \(x\) and \(y\) is a multiple of \((-y,x)\).
	\end{exercise}
	
	\begin{exercise}
		Find two explicit syzygies among
		\[
		x^2,\ xy,\ y^2.
		\]
	\end{exercise}
	
	\begin{exercise}
		Let \(R=k[x,y,z]\). Show that
		\[
		(-y,x,0),\qquad (-z,0,x),\qquad (0,-z,y)
		\]
		are syzygies of \(x,y,z\).
	\end{exercise}
	
	\begin{exercise}
		Let \(M=R/(x,y)\) with \(R=k[x,y]\). Write down a presentation of \(M\) and describe its first syzygy module.
	\end{exercise}
	
	\begin{exercise}
		Explain why \(\Syz(f_1,\dots,f_r)\) is naturally an \(R\)-submodule of \(R^r\).
	\end{exercise}
	
	\section{Solutions sketch}
	
	\subsection*{Exercise 1}
	Compute
	\[
	(-y)x+xy=-xy+xy=0.
	\]
	So \((-y,x)\) is a syzygy.
	
	\subsection*{Exercise 2}
	If
	\[
	ax+by=0,
	\]
	then \(ax=-by\). Since \(x\) and \(y\) are relatively prime, \(y\mid a\). Write \(a=-cy\). Substituting gives \(b=cx\). Hence
	\[
	(a,b)=c(-y,x).
	\]
	
	\subsection*{Exercise 3}
	Two syzygies are
	\[
	(-y,x,0),\qquad (0,-y,x).
	\]
	
	\subsection*{Exercise 4}
	Check directly that each tuple gives zero when substituted into
	\[
	ax+by+cz.
	\]
	
	\subsection*{Exercise 5}
	The quotient \(M=R/(x,y)\) has presentation
	\[
	R^2\to R\to R/(x,y)\to 0,
	\]
	where \(R^2\to R\) sends the standard basis vectors to \(x\) and \(y\). The first syzygy module is generated by \((-y,x)\).
	
	\subsection*{Exercise 6}
	One checks closure under addition and scalar multiplication exactly as in the proof above.
	
	\section{Final summary}
	
	A syzygy is a relation among chosen elements of a ring or module.
	
	The transition from linear algebra to commutative algebra is simple but profound:
	
	\begin{itemize}[leftmargin=2em]
		\item in linear algebra, coefficients come from a field,
		\item in commutative algebra, coefficients come from the ring itself.
	\end{itemize}
	
	Thus, if \(f_1,\dots,f_r\in R\), a syzygy is a tuple
	\[
	(a_1,\dots,a_r)\in R^r
	\]
	such that
	\[
	a_1f_1+\cdots+a_rf_r=0.
	\]
	
	The example
	\[
	(-y)x+xy=0
	\]
	already contains the whole philosophy: generators are not enough by themselves; one must also understand the relations tying them together.
	
	That is what syzygies measure.
	
\end{document}